% ============================================================================
% 00_intro.tex  --  Introduction  (~1/3 page)
% Part 2: Finite-Size Physics & Majorana Edge Modes  (Zhenming Shi)
% Keep the whole report to 4-5 pages; this section is the shortest.
% ============================================================================

\section{Introduction}
\label{sec:intro}

The 1D Kitaev chain is a toy model of a topological superconducting wire,
showing how Majorana zero modes arise at its ends. \textbf{Part 1} studied the
infinite chain with periodic boundaries, where the phase $|\mu| < 2t$ is fixed
by a bulk invariant. A real wire is finite and open, and that finiteness is
what turns ideal zero modes into something measurable: this report treats the
spectrum of the finite open chain, the edge modes and their exponential
localisation, the hybridisation splitting $E_0(L)$, and the sign-filter
artefact that appears once the splitting drops below machine precision. Beyond
the usual sweet-spot ($t = \Delta$) treatment, the analytic backbone is the
exact finite-$L$ theory of Leumer \emph{et
al.}~\cite{leumer2020,leumer2021thesis}: a determinant criterion locating the
$L$ discrete ``Majorana lines'' on which the splitting vanishes \emph{exactly},
and a quantization rule treating bulk and edge states on the same footing.
